%% kapa: exam=1 round=20 year=2009 period=2 form=A subject=appraisal_law subject_ko=부동산관계법규 seq=4 number=44
%% source: Q-Net 감정평가사 기출문제 (cst00309) artlSeq=5024178 file=extracted/1008/5024178/2009년 제20회 감정평가사 1차시험 기출문제/2교시.hwp
%% conversion: kordoc 4.24.51 -> tools/gen1.py (auto); review=proofread-sonnet; flags=tikz
\begin{question}{44}
국토의 계획 및 이용에 관한 법령상 하나의 대지가 아래와 같이 3개의 용도지역에 걸치는 경우 건축할 수 있는 건축물의 지상층 최대 연면적은? (단, 용적률 상한 조건 이외 제약요소 또는 적용요소가 없다고 가정함)
\begin{center}
\begin{tikzpicture}[font=\small, baseline=(current bounding box.center)]
  \draw (0,1.27) rectangle (5.6,5.47);
  \draw (2.43,1.27) -- (2.43,5.47);
  \draw (0,0) rectangle (5.6,1.27);
  \node at (1.2,3.4) {A};
  \node at (4.0,3.4) {B};
  \node at (2.8,0.63) {C};
  \draw (0,-0.7) -- (5.6,-0.7);
  \node at (2.8,-0.4) {도 로};
\end{tikzpicture}\quad
{\small
\renewcommand{\arraystretch}{1.1}
\begin{tabular}{|c|c|c|c|}
\hline
 & 면적 & 용도지역 & \makecell{용적률\\상한} \\
\hline
A & \makecell{300\\제곱미터} & \makecell{제3종\\일반주거지역} & \makecell{250\\퍼센트} \\
\hline
B & \makecell{500\\제곱미터} & 자연녹지지역 & \makecell{80\\퍼센트} \\
\hline
C & \makecell{200\\제곱미터} & 준주거지역 & \makecell{500\\퍼센트} \\
\hline
\end{tabular}}
\end{center}
\begin{choices}
  \item 1,200제곱미터
  \item 1,450제곱미터
  \item 1,650제곱미터
  \item 1,850제곱미터
  \item 2,150제곱미터
\end{choices}
\end{question}
\answer{3,5}
\answernote{복수 정답(정답표 원문: 3, 5)}
